%_ **************************************************************************
%_ * The TeX source for AMS journal articles is the publisher's TeX code    *
%_ * which may contain special commands defined for the AMS production      *
%_ * environment.  Therefore, it may not be possible to process these files *
%_ * through TeX without errors.  To display a typeset version of a journal *
%_ * article easily, we suggest that you either view the HTML version or    *
%_ * retrieve the article in DVI, PostScript, or PDF format.                *
%_ **************************************************************************
%\controldates{1-SEP-1998,1-SEP-1998,1-SEP-1998,13-OCT-1998}

\documentclass{mcom-l} \issueinfo{68}{225}{January}{1999}
\pagespan{389}{393}
\PII{S 0025-5718(99)00991-6}
\def\copyrightyear{1999}

% To avoid problems with the command \MR:
\def\MRhref#1{}

\newcommand{\info}{\(\begin{array}{@{}c@{}}{n\strut}\\{k}\\{\text{digits}}
\\{\text{driver}}\end{array}\)}
\newcommand{\num}[4]{\(\begin{array}{@{}c@{}}{#1\strut}\\{#2}\\{#3}\\{#4}
\end{array}\)}
\newcommand{\pp}{\!\cdot\!}



\begin{document}

\title{Advances in aliquot sequences}

\author{Manuel Benito}
\address{Departamento de Matem\'aticas y Computaci\'on, Universidad de La
Rioja, Edificio J.~L.~Vives, Calle Luis de Ulloa s/n, 26004 Logro\~no, Spain}
\email{mbenito@dmc.unirioja.es}

\author{Juan L. Varona}
\address{Departamento de Matem\'aticas y Computaci\'on, Universidad de La
Rioja, Edificio J.~L.~Vives, Calle Luis de Ulloa s/n, 26004 Logro\~no, Spain}
\email{jvarona@dmc.unirioja.es}
\thanks{Research of the second author supported by grant PB93-0228-C02-02
of the DGICYT}

\subjclass{Primary 11Y55; Secondary 11A25}

\date{October 25, 1996 and, in revised form, February 26, 1997 and July 24,
1997}



\keywords{Aliquot sequences, sum of divisors, perfect number, amicable pair,
sociable numbers, aliquot cycles}

\begin{abstract}
In this paper we describe some advances in the knowledge of the behavior of
aliquot sequences starting with a number less than $10000$.
For some starting values, it is shown for the first time that the sequence
terminates. The current record for the maximum of a terminating sequence is
located in the one starting at~4170; it converges to~1 after~869 iterations
getting a maximum of~84 decimal digits at iteration~289.
\end{abstract}

\maketitle



For a positive integer $n$, let $\sigma(n)$ denote the sum of the divisors of
$n$, and $s(n) = \sigma(n) - n$. A perfect number is a number $n$ such that
$s(n) = n$, and an amicable pair of numbers is $(n,m)$ satisfying
$s(n) = m$, $s(m) = n$. In a similar way, cycles of numbers
$(a_1, a_2, \dots, a_l)$ such that $s(a_i) = a_{i+1}$ for $1 \le i \le l-1$
and $s(a_l) = a_1$ are known as aliquot cycles or sociable numbers.

An aliquot sequence is $\{s^k(n)\}_{k=0}^\infty$ defined by taking
$s^0(n) = n$ and $s^{k+1}(n) = s(s^k(n))$. For one of these sequences,
there are four possibilities:
$(i)$ it terminates at $1$ (the previous term being a prime),
$(ii)$ it reaches a perfect number,
$(iii)$ it reaches an amicable pair or a cycle,
$(iv)$ it is unbounded.
The Catalan-Dickson conjecture~\cite{Cat},
\cite{Dick} says that $(iv)$ does not
actually happen.
But other researchers disagree with this conjecture and think that there
are unbounded sequences; in fact, the alternative conjecture from
Guy-Selfridge~\cite{Guy-Sel} (see also~\cite{Erdos},
\cite{Guy-Guy},
\cite{Guy}) states that
there are many  sequences that go to infinity, perhaps almost all
those that start at an even number.

The nicest way of disproving the Catalan-Dickson conjecture would be to find
an $n$ so that $s^k(n) > s^{k-1}(n)$ for every $k$. Although it seems likely
that such $n$ does not exist, H.~W.~Lenstra has proved that, for every $k$,
there is an $m$ so that $m < s(m) < \cdots < s^k(m)$ (the proof is
constructive and it can be found in~\cite{Erdos}).

Several papers deal with the behavior of these sequences. The reader will
find an extensive bibliography on this subject in~\cite{Guy}. But, before
going on to explain our results, we will summarize the main computational
advances in the study of aliquot sequences.

The smallest $n$ for which there was ever doubt was $138$, but Lehmer showed
that the sequence terminates at $s^{117}(138) = 1$.
Since then, the first number whose behavior is not known is $276$. Many
researchers have investigated the behavior of this sequence (for instance,
G.~A.~Paxson, H.~Cohen, D.~H.~Lehmer, H.~J.~Godwin, J.~L.~Selfridge,
M.~C.~Wunderlich, T.~Struppeck, R.~K.~Guy, A.~Guy and M.~Dickerman),
but it is not yet known if its end exists.

But there are many other sequences whose end is in doubt.
For instance, there are five main sequences starting in a number smaller than
$1000$ whose behavior is unknown (Lehmer five: $276$, $552$, $564$, $660$,
and $966$). Before the year 1980, there were fourteen main sequences starting
between $1000$ and $2000$ (Godwin fourteen), but at the present time only
twelve remain unknown (Godwin showed that the $1848$ sequence terminates, as
did Dickerman for the $1248$ sequence). In~\cite{Guy-Guy}, a table is included
showing the status of the sequences starting with numbers less than or equal
to $7044$. Finally, the latest results in this field can be found in Chapters
B6 and B7 in~\cite{Guy}; here it is stated that ``we have found that those for
$2580$, $2850$, $4488$, $4830$, $6792$, $7752$, $8862$ and $9540$ also
terminate''.


We have been checking the aliquot sequences for numbers under $10000$.
We have used PARI-GP\footnote{C. Batut, D. Bernardi, H. Cohen and M. Olivier,
available by anonymous ftp from \texttt{megrez.ceremab.u-bordeaux.fr} in
\texttt{/pub/pari}.} and, mainly, UBASIC,\footnote{Yuji Kida, available by
anonymous ftp from \texttt{rkmath.rikkyo.ac.jp} in \texttt{/ubibm}.}
and we have run the programs on about twenty computers from the Universidad de
La Rioja and the Instituto P.~M.~Sagasta (Logro\~no, Spain) for about two years,
mostly during nights and weekends. The main factoring method used has been a
combination of both the elliptic curve and the multiple polynomial quad\-ratic
sieve algorithms, and we have checked the primality of the factors with the
APR primality test.
Besides, we wanted to avoid unexpected errors so we have made a further check
of the files containing the sequences. For this task, we have written a
program using the MIRACL\footnote{Michael Scott, available by anonymous ftp
from \texttt{ftp.compapp.dcu.ie} in \texttt{/pub/crypto}.} package.
This program allows us to verify the already computed sequences easily
and quickly; we have executed this program on other computers than the
ones used to calculate the sequences.

For any sequence, we have stopped the calculations when both of the following
happen:
$(i)$ it has a driver (see~\cite{Guy-Sel}), which implies that it is in a
stable increasing situation, without immediate hope of finding its end;
$(ii)$ it has reached a term big enough (75 digits) and whose factorization
is not easily found.
To increase the number of digits of a sequence may require a lot of iterates,
and we are dealing with more than 80 sequences. This involves a very large
amount of computer time.
For instance, we spend about three quarters of an hour to factor a number with
60 decimal digits in a Pentium 100 computer with the multiple polynomial
quadratic sieve algorithm; similarly, one hour and a half for one with
65 digits, 5~hours for one with 70 digits, 18 hours for one with 75 digits, one
day and a half for one with 80 digits, 4~days for one with 85 digits, and
two weeks for one with 90 digits.
So, finally, we stopped computing on February 13, 1997.

In our study, we have achieved several advances over the results of the
aforementioned authors. Some of our conclusions  follow:

\begin{itemize}

\item We have increased the number of iterations in each sequence. In all the
cases, the last term reached has at least $75$ decimal digits.

\item In connection to the most pursued sequence, the one beginning with $276$,
we have reached $913$ iterates without finishing. We have stopped at a term
that has $90$ decimal digits. Also, we have made a special effort with Lehmer
and Godwin sequences.

\item The sequence starting at $8262$ converges to $1$ after $773$
iterations, $317$ being the previous prime. Its maximum is
\[
   \begin{array}{c}
     \hbox{\footnotesize{$
    s^{444}(8262) \,
    = \, 1058123957450778949935095186502241023325216959345026$}}\\
      \hbox{\footnotesize{$
    = \, 2 \cdot 223 \cdot 8831 \cdot 1256153653 \cdot 7873472183
      \cdot 27163313185808805688909099,$}}
   \end{array}
\]
a number of $52$ decimal digits.

\item The sequence starting at $7080$ converges to $1$ after $1264$
iterations, $37$ being the previous prime. Its maximum is
\[
    \begin{array}{c}
      \hbox{\footnotesize{$
    s^{539}(7080) \,
    = \, 3356676040772699493933846784218306089052226182418412008219112234936$}}\\
      \hbox{\footnotesize{$
    = \, 2^3 \cdot 419 \cdot 34110465457 \cdot 1418417145283289
      \cdot 20697305061593509314923291434129565941,$}}
   \end{array}
\]
a number of $67$ decimal digits.

\item The sequence starting at $3556$ converges to $1$ after $2058$
iterations, $59$ being the previous prime. Its maximum is
\smaller[4]
\[
    \begin{array}{c}
      s^{551}(3556) \,
      = \, 476372841845828677729108156155573671704%
        082265291532247964505425538142202028\\
      = \, 2^2 \cdot 503 \cdot 24359 \cdot 7455102267576203135197
      \cdot 100819403542569743 \cdot 12931884982259497162000565321,
    \end{array}
\]\normalsize
a number of $75$ decimal digits.

\item The sequence starting at $4170$ converges to $1$ after $869$
iterations, $79$ being the previous prime. Its maximum is
\smaller[4]
\[
    \begin{array}{c}
     s^{289}(4170) \,
     = \, 329561080342477212747203692863366213833838%
       703158858822327064032192093690321488891836\\
     = \, 2^2 \cdot 41 \cdot 97 \cdot 20374357 \cdot 1559593537
      \cdot 651966073954976081342107597832287652091395156174990523498331163,
    \end{array}
\]\normalsize
a number of $84$ decimal digits. It is a new record for the maximum of a
terminating sequence (the previous record was the $1248$ sequence, which
terminates after reaching a maximum of $58$ digits). We found it on
December 22, 1996.

\item Among the sequences that finish in a cycle, the one that reaches
the largest maximum is the one starting at $7422$. For this sequence, we have
$s^{332}(7422) = 2924$ and $s^{333}(7422) = 2620$, $(2620, 2924)$ being an
amicable pair. The sequence maximum is $s^{156}(7422) = 1269528720481893811316 =
2^2 \cdot 13 \cdot 109 \cdot 131 \cdot 110323 \cdot 15497988349$. (The same
happens with the $7434$ and $9894$ sequences, which coincide with the sequence
corresponding to $7422$ after a few iterations.)

\end{itemize}

Finally, in Table~\ref{tabla} we summarize our work with aliquot sequences
starting with $n < 10000$. In this table we show the main sequences whose status
is yet unknown (by \textit{main} we mean that all the other sequences with
unknown end merge with any sequence in the table; in this case, we only show
in the table the sequence that starts in the smaller number). We also include the
number of decimal digits of the last term $s^k(n)$ reached for each sequence
and the driver in that stage. This table is similar to the one
in~\cite{Guy-Guy}, but now the last term computed for every item is between
$10^{20}$ and $10^{40}$ times bigger (also, the last number analyzed
in~\cite{Guy-Guy} is $7044$). It is important to note that, for every sequence
in the table, we have given up computing iterations at a number containing not
only a guide, but also a driver. As a consequence, no sequences are expected to
start decreasing after only a few iterations.

The complete computations used to work out this table, and the last term
reached for any sequence, can be obtained from the authors upon request. Also,
they are available by anonymous ftp from
\texttt{navalsaz.unirioja.es} in the directory \texttt{/pub/aliquot}.
Equally, we have created a web page about aliquot sequences:
\texttt{http://www.unirioja.es/dptos/dmc/jvarona/aliquot.html}.
This provides a good starting point for future work.


\begin{table}[t]
\caption{Aliquot sequences whose end is in doubt} \label{tabla}
\begin{center}
\renewcommand{\arraystretch}{0.85}
\begin{tabular}{|c|c@{\,}c@{\,}c@{\,}c@{\,}c@{\,}c@{\,}c@{\,}c@{\,}c|}
\hline
\info &
\num{276}{913}{90}{2 \pp 3} &
\num{552}{719}{91}{2^2 \pp 7} &
\num{564}{2230}{91}{2^3 \pp 3} &
\num{660}{333}{90}{2^5 \pp 3 \pp 7} &
\num{966}{358}{91}{2^3 \pp 3 \pp 5} &
\num{1074}{1419}{85}{2^2 \pp 7} &
\num{1134}{2132}{88}{2^6 \pp 127} &
\num{1464}{1703}{88}{2^4 \pp 31} &
\num{1476}{997}{87}{2^3 \pp 3 \pp 5}
\\ \hline
\info &
\num{1488}{580}{84}{2^2 \pp 7} &
\num{1512}{1531}{85}{2^2 \pp 7} &
\num{1560}{1303}{89}{2^5 \pp 3 \pp 7} &
\num{1578}{726}{82}{2^2 \pp 7} &
\num{1632}{653}{82}{2 \pp 3} &
\num{1734}{1262}{89}{2^2 \pp 7} &
\num{1920}{1872}{82}{2^2 \pp 7} &
\num{1992}{949}{88}{2^3 \pp 3 \pp 5} &
\num{2232}{339}{89}{2^3 \pp 3}
\\ \hline
\info &
\num{2340}{410}{78}{2^3 \pp 3 \pp 5} &
\num{2360}{895}{78}{2^2 \pp 7} &
\num{2484}{699}{78}{2^2 \pp 7} &
\num{2514}{2794}{80}{2^3 \pp 3 \pp 5} &
\num{2664}{660}{79}{2^2 \pp 7} &
\num{2712}{800}{82}{2^4 \pp 31} &
\num{2982}{723}{81}{2^4 \pp 31} &
\num{3270}{356}{79}{2^3 \pp 3} &
\num{3366}{988}{79}{2^3 \pp 3}
\\ \hline
\info &
\num{3408}{551}{83}{2^3 \pp 3 \pp 5} &
\num{3432}{882}{85}{2^3 \pp 3 \pp 5} &
\num{3564}{730}{89}{2^3 \pp 3} &
\num{3630}{1030}{80}{2^4 \pp 31} &
\num{3678}{1070}{82}{2^2 \pp 7} &
\num{3774}{1138}{85}{2^3 \pp 3} &
\num{3876}{662}{77}{2^2 \pp 7} &
\num{3906}{638}{83}{2^2 \pp 7} &
\num{4116}{557}{77}{2^3 \pp 3 \pp 5}
\\ \hline
\info &
\num{4224}{459}{78}{2^3 \pp 3 \pp 5} &
\num{4290}{839}{79}{2^2 \pp 7} &
\num{4350}{1062}{80}{2^2 \pp 7} &
\num{4380}{817}{78}{2^2 \pp 7} &
\num{4788}{2095}{82}{2^3 \pp 3 \pp 5} &
\num{4800}{1070}{82}{2^4 \pp 31} &
\num{4842}{365}{79}{2^2 \pp 7} &
\num{5148}{1359}{80}{2 \pp 3} &
\num{5208}{1668}{83}{2^3 \pp 3 \pp 5}
\\ \hline
\info &
\num{5250}{1284}{78}{2 \pp 3} &
\num{5352}{646}{83}{2^2 \pp 7} &
\num{5400}{2077}{76}{2 \pp 3} &
\num{5448}{1154}{86}{2^3 \pp 3 \pp 5} &
\num{5736}{871}{81}{2^2 \pp 7} &
\num{5748}{875}{85}{2^2 \pp 7} &
\num{5778}{680}{78}{2^4 \pp 31} &
\num{6160}{1557}{81}{2^3 \pp 3} &
\num{6396}{1190}{77}{2^3 \pp 3 \pp 5}
\\ \hline
\info &
\num{6552}{845}{82}{2^3 \pp 3} &
\num{6680}{533}{80}{2^3 \pp 3} &
\num{6822}{717}{75}{2 \pp 3} &
\num{6832}{599}{75}{2^2 \pp 7} &
\num{6984}{1568}{75}{2^2 \pp 7} &
\num{7044}{1061}{88}{2^3 \pp 3 \pp 5} &
\num{7392}{463}{83}{2^3 \pp 3 \pp 5} &
\num{7560}{728}{79}{2^4 \pp 31} &
\num{7890}{788}{75}{2^2 \pp 7}
\\ \hline
\info &
\num{7920}{866}{80}{2^6 \pp 127} &
\num{8040}{2172}{83}{2^3 \pp 3 \pp 5} &
\num{8154}{515}{75}{2^2 \pp 7} &
\num{8184}{1122}{79}{2^2 \pp 7} &
\num{8288}{618}{79}{2 \pp 3} &
\num{8352}{1185}{81}{2^2 \pp 7} &
\num{8760}{2091}{75}{2^3 \pp 3 \pp 5} &
\num{8844}{1104}{82}{2^4 \pp 31} &
\num{8904}{834}{80}{2^2 \pp 7}
\\ \hline
\info &
\num{9120}{523}{89}{2^3 \pp 3} &
\num{9282}{373}{77}{2^4 \pp 31} &
\num{9336}{507}{81}{2^6 \pp 127} &
\num{9378}{492}{75}{2^2 \pp 7} &
\num{9436}{316}{75}{2^3 \pp 3} &
\num{9462}{404}{83}{2^3 \pp 3 \pp 5} &
\num{9480}{746}{77}{2 \pp 3} &
\num{9588}{1168}{78}{2 \pp 3} &
\num{9684}{590}{84}{2^5 \pp 3 \pp 7}
\\ \hline
\info &
\num{9708}{558}{83}{2^4 \pp 31} &
\num{9852}{467}{76}{2^2 \pp 7} &&&&&&&
\\ \hline
\end{tabular}
\end{center}
\end{table}

\section*{Acknowledgment} The authors wish to thank the referee for
informing us of the work developed by W.~Creyaufmueller and P.~Zimmermann on
this subject, that appears in the web page
\texttt{http://www.loria.fr/\~{}zimmerma/records/aliquot.html}.
There, it can be seen that they have pursued the sequences starting from
$276$, $552$, $564$, $660$, $966$, $1074$ and $1134$ up to $1059$, $765$,
$2297$, $405$, $423$, $1511$ and $2213$ iterations, respectively (updated to
February 17, 1997). Also, it is announced in the web page that W.~Bosma found
the end of the $3556$ sequence in September 1996.



\bibliographystyle{amsplain}

\begin{thebibliography}{9}

\bibitem{Cat} E. Catalan,
Propositions et questions diverses,
\textit{Bull. Soc. Math. France}
\textbf{18} (1887--88), 128--129.

\bibitem{Dick} L. E. Dickson,
Theorems and tables on the sum of the divisors of a number,
\textit{Quart. J. Math.}
\textbf{44} (1913), 264--296.

\bibitem{Erdos} P. Erd\H{o}s,
On asymptotic properties of aliquot sequences,
\textit{Math. Comp.}
\textbf{30} (1976), 641--645.
\MR{53:7919}

\bibitem{Guy-Guy} A. W. P. Guy and R. K. Guy,
A record aliquot sequence,
in \textit{Computation 1943--1993: A Half-Century of
Computational Mathematics} (Vancouver, 1993),
\textit{Proc. Sympos. Appl. Math.}
\textbf{48} (1994),
Amer. Math. Soc., Providence, RI, 1994,
pp.~557--559.
\MR{95k:11164}

\bibitem{Guy} R. K. Guy,
\textit{Unsolved Problems in Number Theory} (2nd ed.),
Springer-Verlag, 1994.
\MR{96e:11002}

\bibitem{Guy-Sel} R. K. Guy and J. L. Selfridge,
What drives an aliquot sequence?,
\textit{Math. Comp.}
\textbf{29} (1975), 101--107.
Corrigendum, \textit{ibid.} \textbf{34} (1980), 319--321.
\MR{52:5542}; \MR{81f:10008}

\end{thebibliography}

\end{document}


 